\documentclass{abnt}
\usepackage[utf8]{inputenc}
\usepackage[T1]{fontenc}
\usepackage{geometry}
\usepackage{graphicx}
\usepackage{float}
\usepackage{hyperref}
\usepackage[english,indonesian]{babel}

\title{Optimasi Pemrosesan Template LaTeX Menggunakan Pendekatan Berbasis AI}
\author{Academic Support Agent}
\date{\today}

\begin{document}
\maketitle

\begin{abstract}

\end{abstract}

\chapter{PENDAHULUAN}


\chapter{TINJAUAN PUSTAKA}


% Original content (commented out):
% \documentclass[a4paper,12pt,normaltoc,capchap]{abnt}
\usepackage{abnt-UEPG}
\usepackage[top=3cm,bottom=2cm,left=3cm,right=2cm]{geometry}
\usepackage[utf8]{inputenc}
\usepackage[brazil]{babel}
\usepackage[T1]{fontenc}
%\usepackage{hyperref}
\usepackage{amssymb}
\usepackage{amsmath}
\usepackage{paralist}
\usepackage{graphicx}
\usepackage{subfigure}
\usepackage{setspace}
\usepackage{fancyhdr}
\usepackage{lscape}
\usepackage{longtable}
\usepackage{fancyvrb}
\usepackage[alf,recuo=2cm,abnt-etal-text=it,bibjustif]{abntcite}
\usepackage{listings}
\usepackage{fancychap}
\renewcommand{\lstlistingname}{Código} % definição visual dos códigos fonte

\lstset{language=Java, basicstyle=\footnotesize, numbers=left, frame=tb, breaklines=true, basicstyle=\footnotesize}

%------------------------------------------------------------------------------
% Dicionário de hifenização em português
%------------------------------------------------------------------------------

\hyphenation{pro-ces-sa-men-to}
\hyphenation{a-pre-sen-ta-da}
\hyphenation{pro-gra-ma}
%------------------------------------------------------------------------------
% Configurações da dissertação
%------------------------------------------------------------------------------

\autor{UNIVERSIDADE ESTADUAL DE PONTA GROSSA \\
PRÓ-REITORIA DE PESQUISA E PÓS-GRADUAÇÃO\\
PROGRAMA DE PÓS-GRADUAÇÃO EM COMPUTAÇÃO APLICADA\\
CRISTIAN COSMOSKI RANGEL DE ABREU}

\titulo{TÉCNICAS DE COMPUTAÇÃO PARALELA PARA MELHORAR O TEMPO DA MINERAÇÃO DE DADOS:\\Uma análise de Tipos de Coberturas Florestais}

%\orientador{Profº Dr. Luciano José Senger}

%\coorientador{Profº Dr. }

\comentario{Dissertação apresentada para obtenção do título de Mestre em Computação Aplicada na Universidade Estadual de Ponta Grossa, Área de concentração: Computação para Tecnologias Agrícolas. \\ \\Orientador: Prof. Dr. Luciano José Senger}

%\instituicao{UNIVERSIDADE ESTADUAL DE PONTA GROSSA}

\local{Ponta Grossa}

\data{2013}

%------------------------------------------------------------------------------
% Início do documento
%------------------------------------------------------------------------------

\begin{document}

\capa

\folhaderosto

\include{dedicatoria}

\include{agradecimentos}

\include{resumo}

\include{abstract}

%\listadefiguras

%\listadetabelas
\sumario

\pagestyle{headings}

\lhead{\centering \nouppercase{\rightmark}}
\rhead{\centering \nouppercase{\leftmark}}
\pagestyle{fancy}
\fancyfoot{}
\renewcommand{\chaptermark}[1]{\markboth{\thechapter. \ #1}{}}
\fancyhead[LE, RO]{ \thepage}
\fancypagestyle{plain}{
	\fancyhf{}
	\fancyfoot{}
	\renewcommand{\headrulewidth}{0pt}
	\renewcommand{\footrulewidth}{0pt}
}
\pagestyle{fancy}
\pagenumbering{arabic}
\setcounter{page}{1}
\include{introducao}

\include{revisao}

\include{metodologia}

\include{resultados}

\include{conclusoes}

\bibliographystyle{abnt-alf}
\bibliography{base}

%\include{Apendice}

\end{document}

\end{document}